% Standalone research entry 212; batch 208--217.
% Original section material preserved from Pasted text(3).txt.
% Research additions: 27 September 2026.
% Compile twice with: lualatex 212.tex
% !TeX program = lualatex
% Compile twice: lualatex open_problems_500.tex
% All references and prose are embedded; no BibTeX database or external inputs.
\documentclass[11pt,a4paper]{article}
\usepackage[margin=24mm,headheight=14pt]{geometry}
\usepackage{fontspec}
\setmainfont{Latin Modern Roman}
\setsansfont{Latin Modern Sans}
\setmonofont{Latin Modern Mono}
\usepackage{amsmath,amssymb,mathtools}
\usepackage{unicode-math}
\setmathfont{Latin Modern Math}
\usepackage{microtype}
\usepackage{enumitem,xurl,needspace}
\usepackage[dvipsnames]{xcolor}
\usepackage{hyperref}
\hypersetup{unicode=true,colorlinks=true,linkcolor=MidnightBlue,urlcolor=MidnightBlue,
 pdftitle={500 Mathematical Problem Statements},pdfauthor={Source-annotated compilation},bookmarksnumbered=true}
\usepackage{bookmark}
\usepackage{tocloft}
\setlength{\cftsecnumwidth}{3em}
\cftsetpnumwidth{2.5em}
\cftsetrmarg{4em}
\usepackage{titlesec}
\titleformat{\section}{\Large\bfseries\raggedright}{\thesection}{0.7em}{}
\usepackage{fancyhdr}
\pagestyle{fancy}\fancyhf{}
\fancyhead[L]{\small\sffamily Mathematical problem statements}
\fancyhead[R]{\small Source edition 22}
\fancyfoot[C]{\thepage}
\setlength{\parindent}{0pt}
\setlength{\parskip}{5pt plus 1pt}
\setlength{\emergencystretch}{3em}
\setcounter{tocdepth}{1}
\setcounter{secnumdepth}{1}
\setlist[itemize]{leftmargin=3em,itemsep=5pt,topsep=4pt,parsep=0pt}
\allowdisplaybreaks
\newcommand{\N}{\mathbb{N}}
\newcommand{\Z}{\mathbb{Z}}
\newcommand{\Q}{\mathbb{Q}}
\newcommand{\R}{\mathbb{R}}
\newcommand{\C}{\mathbb{C}}
\newcommand{\F}{\mathbb{F}}
\newcommand{\A}{\mathbb{A}}
\newcommand{\PP}{\mathbb P}
\newcommand{\E}{\mathbb{E}}
\newcommand{\eps}{\varepsilon}
\newcommand{\dd}{\,\mathrm d}
\newcommand{\sref}[2]{\hyperlink{source:#1:#2}{[S#2]}}
\newcommand{\eref}[2]{\hyperlink{extra:#1:#2}{[E#2]}}
% Turn the next switch off for a shorter reading copy. The quotations preserve
% every exactTarget field, including qualifications omitted from a short summary.
\newif\ifshowcatalogue
\showcataloguetrue
\newcommand{\cataloguescope}[1]{\ifshowcatalogue
 \paragraph{Exact scope in the supplied catalogue.}
 \begin{quote}\small #1\end{quote}\fi}

% Standalone wrapper; no external inputs or bibliography files are required.
\fancyhead[L]{\small\sffamily Research notebook: section 212}
\fancyhead[R]{\small 27 September 2026}
\hypersetup{pdftitle={Research attempt 212},pdfauthor={Research notebook}}
\setlength{\parskip}{4pt plus 1pt}
\setlist[itemize]{before=\small\interlinepenalty10000,itemsep=3pt}
\begin{document}
\setcounter{section}{211}
% ================================================================
% problemId: problem.morton-silverman-uniform-boundedness-conjecture
% source releaseRank: 212; source releaseStatus: open_with_solved_subcases
% exactTarget SHA-256: e125957f40072ea30fec2c00ba930ac734b7f136ac30223162971602532edabe
\clearpage
\section{\texorpdfstring{Morton–Silverman Uniform Boundedness Conjecture}{Morton–Silverman Uniform Boundedness Conjecture}}\label{212}
\subsection{Definitions and mathematical statement}
A degree-$d$ morphism $\phi:\PP^N_K\to\PP^N_K$ is represented by homogeneous degree-$d$ forms with no common projective zero. A point $P$ is preperiodic when $\{\phi^m(P):m\ge0\}$ is finite, equivalently when $\phi^a(P)=\phi^b(P)$ for some $0\le a<b$. For fixed $N\ge1,d\ge2,D\ge1$, the conjecture is
\[
 \exists C(N,d,D)<\infty\quad
 \#\operatorname{PrePer}(\phi,K)\le C(N,d,D)
\]
for all number fields of degree at most $D$ and every such morphism. The bound cannot depend on the coefficients of $\phi$ or the individual field. Rational maps with points of indeterminacy are outside the selected scope.
\par\smallskip\textit{Formulation and concept references:} \sref{212}{1}.
\cataloguescope{For all fixed integers N>=1, d>=2 and D>=1, there exists a finite constant C(N,d,D) such that for every number field K with [K:Q]<=D and every degree-d K-morphism phi:P\textasciicircum{}N\_K->P\textasciicircum{}N\_K, the set of K-rational preperiodic points has cardinality at most C(N,d,D). A point is preperiodic when its forward orbit under phi is finite. The morphism is everywhere defined, equivalently given by homogeneous degree-d forms with no common projective zero. The bound is independent of the field and map except for N,d,D; no effective formula or classification is required.}
\subsection{Short English statement}
Are the numbers of rational preperiodic points uniformly bounded when projective dimension, map degree, and number-field degree are fixed?
% BEGIN RESEARCH ADDITION 212
\subsection{Research attempt}

\paragraph{Current frontier.}
Northcott finiteness for an individual morphism and number field is not the
coefficient-independent bound requested here; see \sref{212}{1}, Section 4,
Theorem 4.1 and Conjecture 4.2. Looper's revised 2025 manuscript proves a
conditional polynomial result from the $abcd$-conjecture. Its Theorem 1.2 is
stated with a bound $B(d,K)$, so it must not be silently upgraded to a bound
uniform over fields of bounded degree \eref{212}{1}. A July 2026 preprint of
Ji--Xie proves gonality growth in non-isotrivial one-parameter families and
applies it to iterated preimages over number fields and to preperiodic points
over function fields; these are not the unrestricted number-field assertion
\eref{212}{2}, Sections 1.2--1.5. The following local-to-global packing
calculation is a separate attempt, not an invocation of either conditional
arithmetic input.

\paragraph{Attack route.}
There are two distinct quantities to control: the collection of cycles and
the depth of the trees feeding into them. Heights supply finiteness but
normally retain dependence on the map. I instead try to make two distinct
preperiodic points collide in sufficiently small local sets. The product
formula would then forbid the collision. The exact issue is whether the
number of local labels can be bounded independently of the coefficients,
not whether sufficiently fine local covers exist for each fixed map.

\paragraph{Attempt.}
\textit{1. Accounting for tails in the full projective problem.}
A degree-$d$ morphism $\phi:\PP^N\to\PP^N$ is finite of degree $d^N$.
Indeed, a positive-dimensional fiber would contain a projective curve $C$;
$\phi^*\mathcal O(1)|_C=\mathcal O(d)|_C$ would then have both zero and
positive degree, a contradiction. Its degree follows by pulling back the
intersection of $N$ hyperplanes. Every geometric fiber therefore has at most
$d^N$ distinct points, even when it is ramified.

If there are at most $R$ rational periodic points and every rational
preperiodic point reaches a cycle in at most $L$ iterations, it follows that
\[
 \#\operatorname{PrePer}(\phi,K)
 \le R\sum_{j=0}^{L}d^{Nj}.                                    \tag{212.1}
\]
Each point at depth $j$ belongs to the $j$th inverse image of a periodic
point; summing these sets only overcounts. Thus uniform bounds for both
$R$ and $L$ suffice in every dimension, including over fields of degree at
most $D$. A bound just on the number of cycles, with no bound on their
lengths, would not even bound $R$. Neither $R$ nor $L$ has been uniformly
bounded here.

\textit{2. A local covering lemma with an exact global budget.}
Let $K$ be a number field, and normalize absolute values to extend those on
$\Q$, with weights $w_v=[K_v:\Q_v]/[K:\Q]$. Thus
$\prod_v|a|_v^{w_v}=1$ for $a\ne0$. Suppose $X\subset K$ has the following
properties. Outside a finite set $S$, any two elements of $X$ have distance
at most one. At $v\in S$, their distance is at most $R_v>0$. Choose
$T\subseteq S$. For each $v\in T$, cover the image of $X$ in $K_v$ by
$m_v$ sets, each of diameter at most $r_v>0$. If
\[
 \mathcal B:=\prod_{v\in T}r_v^{w_v}
             \prod_{v\in S\setminus T}R_v^{w_v}<1,
 \quad\text{then}\quad |X|\le\prod_{v\in T}m_v.                 \tag{212.2}
\]
To prove this for an arbitrary, possibly infinite $X$, assign a label at
each place by taking the first covering set containing the point. Two
points with the same label tuple would have
\[
 1=\prod_v|x-y|_v^{w_v}\le\mathcal B<1.
\]
Hence the label map is injective, proving (212.2). Diameters, not radii,
are used so that no factor two is lost at archimedean places.

For a monic polynomial $f\in K[z]$, take $X$ to be its finite rational
preperiodic points. They lie in the local filled Julia set, meaning the set
of points with bounded forward orbit. At every non-archimedean place where
all coefficients are integral, $|x|_v>1$ implies
$|f^n(x)|_v=|x|_v^{d^n}$, so $X$ lies in the unit disk. At each remaining
place, elementary domination of the leading term gives a bounded escape
radius. The relevant subsets of $K_v$ are therefore bounded and admit finite
covers of arbitrarily small diameter. This proves the hypotheses preceding
(212.2) with map-dependent data. The point at infinity contributes at most
one additional preperiodic point.

The desired strengthening would be a bound
\[
 \prod_{v\in T}m_v\le C(d,D),\qquad\mathcal B<1,                \tag{212.3}
\]
for every monic degree-$d$ polynomial over every field of degree at most
$D$, with a suitable choice of $T$ and covers. If true, (212.2) would settle
that polynomial subcase without height estimates. It is not asserted that
(212.3) is necessary, and it has not been proved. In particular, replacing
$C(d,D)$ by $d^{|S|}$ leaves the central uniformity problem intact. The role
of local splitting and differences of small points in Looper's proof makes
this obstruction concrete: removing local hypotheses there uses a further
arithmetic conjecture, not the product formula alone \eref{212}{1},
Sections 3--6.

\textit{3. An exact denominator obstruction in the unicritical family.}
Consider $f_c(z)=z^d+c$ and a non-archimedean valuation $v$ on $K$, normalized
with value group $\Z$. If $v(c)<0$ and $x\in K$ is finite and preperiodic,
then
\[
 d\,v(x)=v(c).                                                \tag{212.4}
\]
In absolute-value language, if $|x|_v^d>|c|_v$, its orbit immediately escapes
by strict leading-term domination. If $|x|_v^d<|c|_v$, then
$|f_c(x)|_v=|c|_v>|c|_v^{1/d}$ and the next orbit also escapes. Equality is
the only remaining possibility. This argument does not require the residue
characteristic to be prime to $d$.

For $K=\Q$, (212.4) has an especially sharp consequence. If there is any
finite rational preperiodic point, the reduced denominator of $c$ must be
a perfect $d$th power:
\[
 c=\frac{a}{b^d},\quad (a,b)=1,\quad b\ge1;
 \qquad x=\frac{u}{b},\quad (u,b)=1.                           \tag{212.5}
\]
At primes dividing $b$, (212.4) forces the exact denominator exponent of
$x$. At all other primes, integrality follows from the escape argument.
This is a proved exclusion for many parameters, not just a numerical search.
For example $c=1/A$ with squarefree $A>1$ and $d\ge2$ has no finite rational
preperiodic points; infinity is the only one.

When (212.5) holds, the orbit numerators satisfy
\[
 u_{n+1}=\frac{u_n^d+a}{b^{d-1}}.                              \tag{212.6}
\]
Every surviving rational preperiodic orbit must make this an integer,
coprime to $b$, at every iteration. Thus it obeys a nested congruence sieve,
beginning with $u^d+a\equiv0\pmod{b^{d-1}}$. It also satisfies an archimedean
bound $|u_n|\le b(1+|c|)$: if $|x|>1+|c|$, then
$|f_c(x)|\ge |x|^d-|c|>|x|$ and iteration escapes. For each fixed $c$,
(212.6) therefore gives a finite, exact enumeration procedure by iterating
within the bounded integer interval and rejecting nonintegral transitions.
The size of that interval still depends on $a,b$; this is not a uniform
bound.

\textit{4. Testing whether many bad places actually produce a counterexample.}
Taking $A$ to be the product of more and more primes makes the chosen monic
model $z^d+1/A$ nonintegral at arbitrarily many places, but the previous
lemma makes its finite rational preperiodic set empty. Thus the naive
construction ``more bad places, more preperiodic points'' fails. Conversely,
restricting to parameters whose denominator is a $d$th power avoids this
particular exclusion, but does not automatically give solutions of all the
congruences (212.6). Local admissibility at different primes must be coupled
with the same globally bounded integer orbit.

As small exact checks for $d=2$, $c=-1$ allows just the finite rational
preperiodic points $-1,0,1$, whereas $c=0$ allows the same finite set. For
$c=-1$, every rational preperiodic point is integral, and $|x|\ge2$ escapes;
$-1\mapsto0\mapsto-1$ and $1\mapsto0$. For $c=0$, integrality and real
escape leave $-1,0,1$. Each map has four preperiodic points on $\PP^1(\Q)$
after infinity is included. These examples validate the bookkeeping, not an
extremal conjecture.

\paragraph{Stress test.}
The strict inequality in (212.2) matters: product at most one is consistent
with equality in the product formula and does not prohibit a pair. Covers
at a growing number of places do not have a bounded total label count
merely because each individual count is small. Compactness of one local
filled Julia set gives no uniform covering number across maps. Moreover,
for arbitrary rational functions preperiodic points need not remain in a
fixed bounded affine chart, so the polynomial escape argument cannot simply
be reused. Passing to $\PP^N$ requires additional coordinate and incidence
control, not a change of notation.

The denominator test does not supply the cycle and tail bounds in (212.1).
Nor does a lower bound for the positive canonical heights directly count
height-zero points. Finally, neither a bound depending on a specific field
nor a function-field theorem implies uniformity over all number fields of
bounded degree. These distinctions are retained in the frontier paragraph.

\Needspace{11\baselineskip}
\paragraph{Outcome.}
\textbf{Outcome: Exploratory approach with unresolved gap.}

The attempt proves an adelic collision criterion, an exact denominator and
congruence obstruction for $z^d+c$, and the full-dimensional cycle--tail
counting estimate. It actively tests a bad-reduction construction and shows
why it produces no counterexamples in a broad explicit family. A uniform
local-cover budget is missing even in the proposed monic-polynomial route;
an extension to general projective morphisms is also missing. Nothing here
establishes the unrestricted constant $C(N,d,D)$ or a violating family.
% END RESEARCH ADDITION 212
\subsection{Sources}
\begin{itemize}\raggedright
\item[\textnormal{[S1]}]\hypertarget{source:212:1}{} Current Trends and Open Problems in Arithmetic Dynamics; arXiv1806.04980v2,2July2018; actual PDF received3July2018; journal metadata Bull.AMS56(2019),611–685; HTML internal24August2026 is not a certified scientific revision.
\url{https://arxiv.org/pdf/1806.04980v2}.
{\small\textit{Catalogue formulation source.}}
% BEGIN RESEARCH REFERENCES 212
\item[\textnormal{[E1]}]\hypertarget{extra:212:1}{} Nicole R. Looper,
\emph{The Uniform Boundedness and Dynamical Lang Conjectures for polynomials},
arXiv:2105.05240v3 (21 December 2025), Theorem 1.2, Conjecture 2.1, and
Sections 3--6. The stated bound in Theorem 1.2 is $B(d,K)$.
\url{https://arxiv.org/html/2105.05240v3}.
\item[\textnormal{[E2]}]\hypertarget{extra:212:2}{} Zhuchao Ji and Junyi Xie,
\emph{Genus and Gonality of Small Curves, Dynamical Uniform Boundedness, and Bifurcation},
arXiv:2607.12561v2 (18 July 2026), Sections 1.2--1.5 and 8.
\url{https://arxiv.org/html/2607.12561v2}.
{\small\textit{Recent preprint; the applications cited are not the unrestricted number-field conjecture.}}
% END RESEARCH REFERENCES 212
\end{itemize}

\end{document}
